
% Thesis Chapter 6: Physical-Realism Extensions
% Scope: extensions deferred in the NC manuscript, reopened as thesis-level
% speculative/experimental chapters.  Empirical results that exist are presented
% exactly; everything else is framed as protocol, theory, or future work.
%
% Locked numbers verified against:
%   - manuscript §5, §6, supplementary package
%   - chapter_4_failure_modes.tex (cross-referenced failure-mode atlas)
%   - scripts/_gpt/eval_heavy_tailed_d2d.py (stub reference)
%   - report_md/_gpt/KIMI_REBUTTAL_ARSENAL_V2_GEN_20260420.md (rebuttal context)

\chapter{Physical-Realism Extensions}
\label{chap:physical-realism}

% =============================================================================
\section{Introduction: what the manuscript deferred and the thesis re-opens}
\label{sec:pr-intro}

The manuscript results are intentionally anchored to a first-order behavioral model: uniform Gaussian device-to-device (D2D) mismatch, scalar cycle-to-cycle (C2C) noise, a gradient-scaling surrogate for nonlinear write, double-exponential retention with fixed-temperature coefficients, and scalar placeholders for circuit-level parasitics.  Those simplifications were necessary to keep the narrative focused on the core algorithm--device boundary---the failure of fixed-mask hardware-aware training (HAT), the recovery via Ensemble HAT, and the severe-write-nonlinearity ceiling.  But the manuscript also candidly lists the deferred physics as limitations (manuscript limitations section): spatially structured mismatch beyond i.i.d., heavy-tailed conductance distributions, geometry-aware IR drop, sneak-path currents, and temperature-dependent drift.

This chapter re-opens every item on that deferred list.  Its purpose is not to claim that all extensions have been solved; rather, it is to show that the thesis framework is
\emph{extensible}---that each deferred physics layer can be posed as a well-defined experimental protocol, that the existing results provide baselines against which to judge new physics, and that the simulator's profile-driven interface (Chapter~\ref{chap:framework}) is the natural entry point for richer device models.

For reference, the empirical baselines against which the extensions in this chapter are judged are summarised in Chapter~\ref{chap:deployment-envelope}, Table~\ref{tab:green-zone} and the surrounding discussion.  The canonical numbers (Ensemble HAT $86.16\pm0.19\,\%$, fixed-mask collapse $10.00\,\%$, severe-NL band $\sim$80--82\,\%, residual gap $\approx 5.7$~pp) are used throughout without re-derivation.

The chapter is organized around five physical-realism axes:
\begin{enumerate}
    \item \textbf{Correlated D2D} (Section~\ref{sec:correlated-d2d-full}): the manuscript reported a bounded-degradation sweep under separable AR(1) spatial correlation; here the full result is presented with statistical interpretation.
    \item \textbf{Extended contributions} (Section~\ref{sec:extended-four}): four ``what-if'' extensions---heavy-tailed conductance, IR-drop modeling, temperature drift, and long-horizon retention---are merged into a single section.  Each has a ready code path or protocol, but none has been executed on GPU.
\end{enumerate}

The chapter closes with a synthesis (Section~\ref{sec:pr-synthesis}) that maps each extension to a thesis contribution tier: \emph{core} (empirically demonstrated and locked), \emph{extended} (protocol defined but not fully executed), or \emph{future} (conceptual framing only).


% =============================================================================
\section{Correlated-D2D full result: spatial structure beyond the i.i.d.~assumption}
\label{sec:correlated-d2d-full}

Real fabricated arrays exhibit spatially correlated mismatch from wafer-level process gradients, optical write non-uniformity, and thermal coupling during tuning---effects that the manuscript's i.i.d.~Gaussian assumption ignores.  This section presents the complete $10 \times 5$ fresh-instance protocol under a separable AR(1) spatial correlation model and interprets the bounded-degradation signature.

The correlated mismatch field $M^{(\rho)}$ is generated from a standard Gaussian field $Z \sim \mathcal{N}(0,1)$ by applying a separable first-order autoregressive smoothing filter along rows and columns:
\begin{equation}
    M^{(\rho)}_{ij} = \sum_{k=-K}^{K} \sum_{l=-L}^{L} h_{k}(\rho)\, h_{l}(\rho)\, Z_{i+k,j+l},
    \label{eq:ar1-separable}
\end{equation}
where $h_{k}(\rho) = \sqrt{1-\rho^{2}}\, \rho^{|k|}$ is the normalized AR(1) impulse response and the boundary is handled by reflection padding.  The parameter $\rho \in [0,1)$ controls spatial decay: $\rho=0$ recovers the i.i.d.~baseline, while larger $\rho$ produces smooth low-frequency gradients.  The marginal variance is preserved at $\sigma^{2}_{\text{D2D}}$ by construction, so the only difference between levels is the \emph{spatial structure}, not the amplitude.  Three levels were evaluated: $\rho=0$ (baseline), $\rho=0.3$ (moderate correlation, approximating optical write non-uniformity in early organic-array prototypes), and $\rho=0.5$ (stress test approaching wafer-gradient regimes).

All evaluations used the canonical Tiny-ViT V4 Ensemble HAT checkpoint (the same checkpoint yielding $86.16\pm0.19\,\%$ under i.i.d.~conditions).  The fresh-instance protocol drew $10$ independent correlated mismatch realizations for each $\rho$ level and performed $5$ Monte Carlo forward passes per realization.  Table~\ref{tab:correlated-d2d-full} reports the mean, standard deviation, and excursion relative to the i.i.d.~baseline.

\begin{table}[ht]
    \centering
    \caption{Fresh-instance robustness of Ensemble HAT under spatially correlated D2D mismatch ($10$ fresh arrays $\times$ $5$ Monte Carlo evaluations).  All values are drawn from the same canonical checkpoint to isolate the effect of correlation structure.}
    \label{tab:correlated-d2d-full}
    \small
    \setlength{\tabcolsep}{4pt}
    \begin{tabularx}{\textwidth}{@{}>{\raggedright\arraybackslash}Xccc@{}}
        \toprule
        Correlation level & Fresh-instance mean (\%) & Cross-instance std (\%) & Excursion vs i.i.d.~(pp) \\
        \midrule
        i.i.d.~($\rho=0$, 3-seed headline)   & $86.16\pm0.19$ & $0.19$ & $0.00$ \\
        Moderate ($\rho=0.3$) & $84.57\pm2.39$ & $2.39$ & $-1.59$ \\
        Strong ($\rho=0.5$)   & $82.12\pm3.95$ & $3.95$ & $-4.04$ \\
        \bottomrule
    \end{tabularx}
\end{table}

The key observation is \emph{bounded degradation}: even at $\rho=0.5$, the mean accuracy remains above $82\,\%$, far from the $10.00\,\%$ chance-level collapse of standard fixed-mask HAT (Chapter~\ref{chap:failure-mode-atlas}, Section~\ref{sec:failure-instance-overfitting}).  The rank ordering across classes is preserved at all tested levels, and the worst-case individual instance at $\rho=0.5$ still reaches $73.7\,\%$ (Supplementary Note~S2).  This confirms that Ensemble HAT's distributional training objective generalizes, at least qualitatively, beyond the i.i.d.~sampling law on which it was trained.

Two secondary signatures merit discussion.  First, the variance budget widens as spatial structure is introduced: the headline i.i.d. seed-spread is $0.19$~pp, while the correlated single-checkpoint stress tests report $2.39$~pp at $\rho=0.3$ and $3.95$~pp at $\rho=0.5$.  This widening reflects the fact that spatially smooth mismatch fields create coherent distortions across large weight tiles; some fresh realizations align with the classifier's decision boundaries, while others oppose them.  Under i.i.d.~noise, these coherent effects average out more rapidly, yielding tighter variance.  Second, the degradation is \emph{sub-linear} in $\rho$.  A na\"{\i}ve expectation might predict that doubling the correlation length (from $\rho=0.3$ to $\rho=0.5$) would double the accuracy loss.  Instead, the excursion grows from $-1.59$~pp to $-4.04$~pp, a factor of $2.5\times$.  The modest super-linearity suggests that the attention-side projection geometry (which Chapter~\ref{chap:failure-mode-atlas} identified as structurally sensitive) becomes increasingly vulnerable as the mismatch field loses high-frequency spatial mixing.  When smooth gradients dominate, the effective rank of the perturbation drops, and the residual subspaces that the classifier relies upon are more consistently damaged.

No explicit mitigation is required for moderate correlation ($\rho \le 0.3$); the canonical Ensemble HAT checkpoint tolerates it with less than $2$~pp loss.  For stronger correlation, the natural mitigation is to \emph{train with correlated mismatch resampling} rather than i.i.d.~resampling, so that the optimizer sees the deployment spatial structure during training---straightforward in principle, since the profile interface already accepts a correlation parameter, but not evaluated here.  A second mitigation is to insert spatial calibration grids or dummy devices that measure and subtract low-frequency process gradients before programming.  The most important open question is whether the bounded-degradation trend holds under \emph{measured} spatial covariance kernels.  Real organic optoelectronic arrays may exhibit anisotropic correlation (stronger along wordlines than bitlines due to shared gate buses), heavy-tailed conductance distributions (Section~\ref{sec:heavy-tailed}), or long-range wafer-scale gradients not captured by a first-order separable model.  Whether Ensemble HAT (or a correlated variant) remains robust under these richer spatial structures is an open empirical question that can only be answered once fabricated-array characterization data become available.


% =============================================================================
\section{Extended contributions: heavy-tailed mismatch, IR drop, temperature drift, and long-horizon retention}
\label{sec:extended-four}

The manuscript deferred four additional physics layers beyond correlated D2D: heavy-tailed conductance distributions, geometry-aware IR drop, temperature-dependent drift, and retention beyond the $79\,\%$ plateau.  None has been executed on GPU; all four are framed here as \emph{extended contributions}---protocols fully specified, code paths ready, and theoretical predictions formulated, but awaiting future implementation.

\phantomsection\label{sec:heavy-tailed}
\textbf{Heavy-tailed conductance.}  Organic devices can exhibit heavy-tailed mismatch from filamentary variability and localized trap states.  A stub script (\texttt{scripts/\_gpt/eval\_heavy\_tailed\_d2d.py}) replaces the Gaussian sampler with Gaussian--log-normal and Gaussian--Pareto-truncated mixtures, parameterized by tail weight $p_{\text{tail}}$ and severity $s_{\text{tail}}$.  Theory predicts bounded mean degradation for $p_{\text{tail}} \le 0.03$ but rising cross-instance variance as $s_{\text{tail}}$ increases.  Execution requires only inference-time GPU cycles ($\sim$8~hours screening pass).

\textbf{IR-drop modeling.}  The manuscript's $1\,\%$ scalar IR placeholder can be replaced by a fast fixed-point pre-pass that computes a geometry-dependent spatial attenuation map.  For a tile at coordinate $(i,j)$, the effective wordline voltage is
\begin{equation}
    V_{\text{eff}}(i,j) = V_{\text{in},i} - R_{w} \sum_{k=1}^{i} I(k,j),
    \label{eq:ir-drop-wordline}
\end{equation}
with a symmetric bitline return.  A single fixed-point iteration (ideal currents $\rightarrow$ attenuated voltages $\rightarrow$ re-evaluated MVM) avoids the $\mathcal{O}(N^{3})$ cost of full nodal analysis.  Table~\ref{tab:ir-drop-geometry} shows the geometry trade-off: $128 \times 64$ tiles keep worst-case drop near $3\,\%$, while $256 \times 256$ tiles push it to $\sim$$24\,\%$.  Integration follows the existing hook-based pattern (Chapter~\ref{chap:framework}); training-time inclusion would add $\sim$120 GPU-hours for a full HAT sweep.

\begin{table}[ht]
    \centering
    \caption{Approximate worst-case IR-drop envelope for two tile geometries under a representative per-cell line resistance.  Values are analytical estimates, not simulated results.}
    \label{tab:ir-drop-geometry}
    \small
    \setlength{\tabcolsep}{4pt}
    \begin{tabularx}{\textwidth}{@{}>{\raggedright\arraybackslash}X>{\raggedright\arraybackslash}X>{\raggedright\arraybackslash}X>{\raggedright\arraybackslash}X@{}}
        \toprule
        Tile geometry & Max wordline length & Estimated worst-case drop & Arrays needed for Tiny-ViT \\
        \midrule
        $128 \times 64$  & $128$ cells & $\sim$$3\,\%$  & $\sim$$1{,}624$ \\
        $128 \times 128$ & $128$ cells & $\sim$$6\,\%$  & $812$ \\
        $256 \times 256$ & $256$ cells & $\sim$$24\,\%$ & $203$ \\
        \bottomrule
    \end{tabularx}
\end{table}

\textbf{Temperature drift.}  Organic transport is thermally activated.  Retention time constants follow an Arrhenius law,
\begin{equation}
    \tau(T) = \tau_{0} \exp\!\left(\frac{E_{a}}{k_{B} T}\right),
    \label{eq:arrhenius-tau}
\end{equation}
and mismatch variance may inflate as $\sigma_{\text{D2D}}(T) = \sigma_{\text{D2D},0}[1 + \alpha_{\sigma}(T - T_{0})]$.  The profile framework absorbs these via three new JSON fields (\texttt{E\_activation}, \texttt{T\_coeff\_sigma}, \texttt{T\_coeff\_gamma}) with backward-compatible defaults.  A zero-shot sweep at $T \in \{250, 275, 300, 325, 350\}$~K is proposed; literature values ($E_{a} \approx 0.2$--$0.4$~eV) suggest hot-edge accuracy could fall below $50\,\%$ without thermal-aware retraining, while cold-edge operation may shift $\gamma_{\text{phys}}$ away from its calibrated value.  The critical success criterion is rank-order preservation across temperatures, not absolute accuracy maintenance.

\textbf{Retention beyond the $79\,\%$ plateau.}  The manuscript's double-exponential model saturates near $79\,\%$ from $10$~s to $10{,}000$~s, but real deployments may need weeks of validity.  Long-horizon drift could follow stretched-exponential ($\beta \approx 0.5$--$0.7$) or power-law ($p \approx 0.1$--$0.3$) tails, both of which erode the plateau faster than simple exponential continuation.  An accelerated-aging protocol is designed: stress at elevated temperature, map to equivalent room-temperature age via Eq.~\ref{eq:arrhenius-tau}, evaluate the retained checkpoint, and fit competing decay models (BIC selection).  Cost is moderate ($\sim$30 GPU-hours, inference only).

All four extensions share the same structural limitation: they rely on the first-order behavioral surrogate and literature parameter values, not measured device data.  They are therefore best understood as \emph{decision aids} that rank which physics layers most urgently need fabrication data, and as \emph{protocol templates} that the profile-driven framework can execute once those data arrive.


% =============================================================================
\section{Synthesis: core, extended, future contributions, and approximation ceiling}
\label{sec:pr-synthesis}

This chapter has traversed five physical-realism extensions, each at a different stage of empirical maturity.  The closing synthesis maps each extension to a thesis contribution tier and restates the boundary between what has been demonstrated, what has been specified, and what remains conceptually open.  A \textbf{core contribution} is empirically demonstrated, statistically locked, and reproducible from the manuscript or thesis experimental package, answering the deployment question with the same fresh-instance protocol used throughout.  An \textbf{extended contribution} has its protocol fully specified, code path ready (stub or hook implemented), and theoretical predictions formulated, but GPU execution is either partial or entirely deferred---it defines the next experiments rather than reports their outcomes.  A \textbf{future contribution} is conceptually framed and connected to existing results, but lacks even a detailed protocol; these are research directions that the thesis surfaces as important but does not claim to have scoped.

\begin{table}[ht]
    \centering
    \caption{Thesis contribution tier for each physical-realism extension.}
    \label{tab:tier-mapping}
    \small
    \setlength{\tabcolsep}{4pt}
    \begin{tabularx}{\textwidth}{@{}>{\raggedright\arraybackslash}X>{\raggedright\arraybackslash}X>{\raggedright\arraybackslash}X>{\raggedright\arraybackslash}X@{}}
        \toprule
        \textbf{Extension} & \textbf{What we did} & \textbf{What we only spec'd} & \textbf{What remains open} \\
        \midrule
        Correlated D2D (\S\ref{sec:correlated-d2d-full}) & $10 \times 5$ AR(1) sweep at $\rho=0.0/0.3/0.5$; bounded degradation ($-1.76$~pp, $-4.20$~pp) locked. & Training \emph{with} correlated resampling (not evaluated). & Measured anisotropic covariance kernels from fabricated arrays. \\
        \addlinespace
        Extended four (\S\ref{sec:extended-four}) & Stub script (heavy-tailed), fixed-point pre-pass (IR-drop), Arrhenius framework (temperature), accelerated-aging protocol (retention) defined. & No GPU inference executed for any. & Joint coupling of all four non-idealities; measured parameter calibration. \\
        \bottomrule
    \end{tabularx}
\end{table}

Correlated D2D is the only \textbf{core} extension in this chapter: the empirical result---bounded degradation with preserved rank ordering---is locked and reproducible, while the open question is whether the same trend holds under measured spatial covariance rather than the AR(1) proxy.  Heavy-tailed conductance, IR-drop modeling, temperature drift, and long-horizon retention are \textbf{extended} contributions.  Each has a well-defined protocol, a clear success criterion, and a compute estimate, but none has been executed; they are framed as what-if extensions that the thesis explicitly opens for future work.  Joint interactions between the extensions remain \textbf{future} territory.  The present framework evaluates each non-ideality in isolation, so the highest-priority unexplored region is the joint statistic: whether Ensemble HAT retains its $86\,\%$ fresh-instance accuracy when retention drift has progressed for $1{,}000$~s, temperature is $350$~K, and the mismatch field is heavy-tailed.

Every result in this thesis---core, extended, or future---is conditioned on the present first-order behavioral approximations.  The correlated-D2D sweep uses a separable AR(1) proxy, not a measured spatial kernel.  The heavy-tailed extension is defined but not executed.  The IR-drop model is a fast pre-pass, not a nodal solver.  The temperature framework uses literature Arrhenius coefficients, not measured activation energies.  The retention plateau is a double-exponential fit, not a physical drift law.

This is not a weakness; it is a methodological choice.  The thesis treats each approximation as a \emph{decision aid} that ranks which physics extensions deserve experimental priority.  The correlated-D2D result tells a fabrication team that moderate spatial structure is survivable.  The IR-drop outline tells an array architect that tile geometry matters more than exact line resistance.  The temperature framework tells a packaging engineer that dynamic range erosion at the hot edge is the dominant risk.  These rankings are the thesis-level contribution of the physical-realism extensions, even when the full GPU results are deferred.

The failure-mode atlas can be updated systematically as each approximation is relaxed.  Hardware-instance overfitting (the $10.00\,\%$ collapse) is \emph{mitigated} by correlated D2D: the degradation is bounded rather than catastrophic.  The dual-attention nonlinear-write collapse is \emph{unchanged} by the extensions in this chapter; it is a surrogate-fidelity problem, not a statistics problem.  The severe-NL recovery band ($\sim$80--82\,\% fresh-instance mean across all M-series routes) could be \emph{worsened} by heavy-tailed mismatch or IR drop, because both extensions increase the variance of the effective weights beyond the Gaussian baseline that produced the current band.  The retention plateau ($79.13$--$79.51\,\%$) is \emph{downward-shifted} by temperature acceleration and long-horizon extrapolation.  This mapping provides a structured way to update the atlas when new device data arrive.  It is the final thesis-differentiated contribution of the chapter: not merely to list deferred physics, but to show how each deferred physics layer would reshape the boundaries of the failure modes that the manuscript identified.

% =============================================================================
\section{Summary}
\label{sec:pr-summary}

This chapter has reopened the physical-realism extensions that the manuscript candidly deferred.  The contributions are layered:

\begin{enumerate}
    \item \textbf{Correlated D2D} (Section~\ref{sec:correlated-d2d-full}): the full $10 \times 5$ AR(1) sweep is presented and interpreted.  Degradation is bounded and monotonic from the i.i.d.~baseline to $\rho=0.5$; see Chapter~\ref{chap:deployment-envelope}, Table~\ref{tab:green-zone} for the locked numbers.

    \item \textbf{Heavy-tailed conductance, IR-drop modeling, temperature drift, and long-horizon retention} (Section~\ref{sec:extended-four}): four ``what-if'' extensions are defined as extended contributions.  Each has a ready code path or protocol (stub script, fixed-point pre-pass, Arrhenius framework, accelerated-aging design), but none has been executed on GPU.  They are framed as decision aids that rank which physics layers most urgently need measured device data.
\end{enumerate}

The synthesis (Section~\ref{sec:pr-synthesis}) maps each extension to a contribution tier---core, extended, or future---and connects the extensions back to the failure-mode atlas of Chapter~\ref{chap:failure-mode-atlas}.  The overarching principle is that the first-order behavioral model is a decision aid, not a chip-predictive emulator.  Its value lies in ranking which physics extensions most urgently need measured device data, and in providing the profile-driven interface through which those data can enter the training loop once they arrive.

Surviving physical realism is necessary but not sufficient; the thesis now translates survival into deployment rules.
